\documentclass[11pt,a4paper]{article}
\usepackage{iftex}
\ifXeTeX
  \usepackage{fontspec}
  \setmainfont{Liberation Serif}
  \setsansfont{Liberation Sans}
  \setmonofont{Liberation Mono}
\else
  \usepackage[T2A]{fontenc}
  \usepackage[utf8]{inputenc}
\fi
\usepackage[russian]{babel}
\usepackage[margin=2cm]{geometry}
\usepackage{amsmath,booktabs,tabularx,array}
\usepackage{graphicx}
\usepackage[unicode,hidelinks]{hyperref}
\usepackage{xurl}
\setlength{\parindent}{1.25cm}
\setlength{\parskip}{3pt}
\setlength{\emergencystretch}{2em}
\renewcommand{\arraystretch}{1.15}
\renewcommand{\tablename}{Таблица}
\title{Обзор отраслевого отчёта\\
«Топ-40 стран-импортёров мира. ТН ВЭД 040900\\
«Мёд натуральный», объём в тоннах»\\
(платформа GTAIC / EXPORT HUNTER, 2026)}
\author{Команда 10130\\Курс «Поиск экспортных рынков»}
\date{25 ноября 2026 г.}
\begin{document}
\maketitle

\noindent Настоящий обзор подготовлен в рамках семинарской активности курса. Команда исследует
экспортные возможности России на рынке натурального мёда (HS6 040900) и сравнивает условия
поставки на трёх целевых рынках -- США, Японии и Южной Корее. Предмет обзора -- отраслевой
отчёт платформы GTAIC (издатель EXPORT HUNTER, UAB, Вильнюс)~\cite{gtaic2026}: документ объёмом
217~страниц, сгенерирован 24.11.2026 на данных UN Comtrade, Eurostat (Comext) и китайской
таможни GACC; охват -- топ-40 стран-импортёров, на которые приходится 96{,}07\,\% мирового
импорта мёда в тоннах. Нам в этом отчёте важнее всего то, чего нет в нашей долларовой панели
BACI, -- физические объёмы в тоннах и средние цены-прокси. Отчёт также ставит вопрос
о масштабе и динамике спроса на целевых рынках и даёт контекст для интерпретации: позиции
конкурентов, торговые меры, валютный фон. Собственные расчёты команды выполнены
по OEC/CEPII BACI в долларах США за 2018--2024~\cite{oecbaci}; обеспеченность пашней --
из WDI/FAO~\cite{wdi}; товарная характеристика мёда -- по FAO~\cite{bradbear}. Ниже -- паспорт и
методология отчёта (разд.~1), общая динамика рынка (разд.~2), крупнейшие импортёры и наши
целевые рынки (разд.~3), поставщики и ценовая структура (разд.~4), сопоставление отчёта с
нашими расчётами (разд.~5), критическая оценка (разд.~6) и значение для дальнейшей работы
(разд.~7).

\section*{1. Паспорт отчёта и методология}

Состав отчёта: паспорт продукта и методология (стр.~1--13), сводная статистика и рейтинги
(стр.~15--175), постраничные профили 40 стран-импортёров (стр.~177--216). Базовый период --
календарные 2020--2025 годы; текущие данные приведены на двенадцатимесячном интервале LTM
(last twelve months), у каждой страны своё окно (стр.~12).

\begin{table}[h]
\centering
\small
\caption{Паспорт отчёта GTAIC по рынку мёда (HS 040900)}
\begin{tabular}{p{4.2cm}p{11.2cm}}
\toprule
Издатель & EXPORT HUNTER, UAB, Вильнюс; платформа www.gtaic.ai \\
Дата генерации & 24 ноября 2026 г.; объём 217 стр. \\
Товар & HS 040900 «Мёд натуральный» \\
Охват & топ-40 импортёров -- 96{,}07\,\% мирового импорта в тоннах (LTM) \\
Базы данных & UN Comtrade (свыше 200 стран); Eurostat Comext -- импорт Австрии и Франции,
экспортная статистика Австрии, Кипра, Франции, Турции; GACC (Китай) -- с 01.2025 \\
Периоды & годовые 2020--2025; частичный 2026*; LTM-окна по каждой стране отдельно \\
Измерители & стоимость (млн долл.), тоннаж (тыс.~т), средняя цена-прокси CIF (долл./т) \\
\bottomrule
\end{tabular}

\noindent Источник: GTAIC (2026), обложка и стр.~1--13.
\end{table}

В методологии четыре особенности, существенные для чтения цифр. Во-первых, LTM-окна
не синхронизированы: у большинства стран это сентябрь 2025 -- август 2026, у Японии и ещё
пяти стран -- октябрь 2025 -- сентябрь 2026, у Республики Корея*, ОАЭ*, Кувейта*, Тайваня*
и Марокко* -- май 2025 -- апрель 2026, у Сингапура -- апрель 2025 -- март 2026 (стр.~12).
Суммарный «агрегат LTM» поэтому складывает разные периоды, и межстрановые сравнения внутри
LTM-разделов формально некорректны; корректнее опираться на календарные 2020--2025 годы.
Во-вторых, для пяти стран (Тайвань*, Кувейт*, Марокко*, Республика Корея*, ОАЭ*) импорт
реконструируется «зеркально» по экспортным данным 62 стран-партнёров, и отчёт прямо пишет,
что такие оценки «могут быть занижены»: непокрытыми остаются, в частности, 8{,}66\,\%
импорта Кореи* (Сингапур) и 5{,}28\,\% (Сингапур и др.). В-третьих, если источник не сообщает
вес, тоннаж пересчитывается по количеству мест или по медианным ценам с градациями
достоверности; доля заявленного веса колеблется от 76{,}87\,\% у США и 78{,}84\,\% у Германии
до 0\,\% у Израиля, а весь вес Китая с 2025 года расчётный, поскольку GACC публикует только
стоимости. В-четвёртых, отчёт официально предупреждает, что «весьма вероятно», Китай,
Республика Корея* и ЮАР сообщили неверные объёмы в тоннах за LTM из-за аномалий цен CIF
(в два раза выше или ниже медианных), и рекомендует проверять тоннаж этих трёх стран по другим
источникам. Последнее замечание для нашей команды критично, поскольку Южная Корея входит
в число целевых рынков.

\section*{2. Общая картина рынка: рекорд по объёму, но не по деньгам}

Календарные ряды рынка показаны в табл.~2. Главный результат 2025 года -- исторический максимум
физического импорта топ-40: 776{,}05 тыс.~т (+2{,}1\,\%) при 2{,}32 млрд долл. (+8{,}0\,\%).
В деньгах рынок остаётся примерно на 11\,\% ниже пика 2022 года (2{,}62 млрд долл. при
726 тыс.~т). Пятилетний CAGR составляет $+2{,}02\,\%$ по тоннажу и $+1{,}60\,\%$ по стоимости,
а средняя цена-прокси -- $-0{,}42\,\%$ в год. Частичный 2026* (одни и те же месяцы соседних
лет) продолжает восстановление: $+5{,}75\,\%$ тонн и $+9{,}13\,\%$ долларов.

\begin{table}[h]
\centering
\small
\caption{Импорт мёда странами топ-40 по годам, 2020--2026*}
\begin{tabular}{lrrr}
\toprule
Период & Импорт, млрд долл. & Тоннаж, тыс.~т & Средняя цена, долл./т \\
\midrule
2020 & 2{,}14 & 702 & 3\,050 \\
2021 & 2{,}53 & 707 & 3\,590 \\
2022 & 2{,}62 (пик) & 726 & 3\,610 (пик) \\
2023 & 2{,}09 & 650 (мин.) & 3\,210 \\
2024 & 2{,}14 & 760 & 2\,820 (мин.) \\
2025 & 2{,}32 & 776 (макс.) & 2\,983 \\
2026* (част.) & 1{,}55 & 518 & 2\,989 \\
\bottomrule
\end{tabular}

\noindent Источник: GTAIC (2026), рис.~14--16. Единицы: млрд долл. США, тыс.~т, долл./т;
2026* -- доступные месяцы 2025 и 2026 гг.\ (сопоставимый частичный период), без LTM-агрегатов.
\end{table}

Последовательность лет объяснима ценовой динамикой: пиковые цены 2021--2022 годов
(до 3\,610 долл./т) совпали с провалом тоннажа 2023 года до 650 тыс.~т, после чего объёмы
два года восстанавливались на фоне удешевления. Мы фиксируем отрицательную связь цены
и физического объёма на уровне рынка; говорить о причинной интерпретации (эластичности
спроса) по этим данным нельзя -- не учтены сортовая структура поставок, доходы и курсы валют.
Для нашего исследования важно то, что рынок растёт именно в физическом измерении: спрос
восстанавливается преимущественно в массовых ценовых сегментах.

\section*{3. Крупнейшие импортёры и наши целевые рынки}

Сводка по десяти крупнейшим рынкам LTM приведена в табл.~3; там же для сопоставления -- Южная
Корея (24-е место). 27 из 40 рынков в отчёте -- европейские. На пятёрку лидеров приходится
58{,}0\,\% стоимости 2025 года.

\begin{table}[h]
\centering
\small
\caption{Топ-10 стран-импортёров мёда, LTM (и Южная Корея)}
\begin{tabular}{rlrrrrr}
\toprule
№ & Страна & млн долл. & $\Delta$\,\%, \$ & тыс.~т & $\Delta$\,\%, т & CAGR 5 лет, \% \\
\midrule
1 & США & 736{,}33 & $+4{,}1$ & 264{,}11 & $+1{,}3$ & $+9{,}36$ \\
2 & Германия & 263{,}78 & $+10{,}3$ & 73{,}98 & $-0{,}7$ & $-2{,}03$ \\
3 & Япония & 142{,}32 & $-3{,}4$ & 48{,}32 & $-5{,}7$ & $-2{,}84$ \\
4 & Великобритания & 141{,}64 & $+10{,}0$ & 62{,}40 & $+3{,}8$ & $+1{,}95$ \\
5 & Франция & 127{,}05 & $+20{,}1$ & 33{,}02 & $+9{,}7$ & $-2{,}84$ \\
6 & Испания & 98{,}86 & $+53{,}6$ & 44{,}35 & $+47{,}6$ & $+3{,}04$ \\
7 & Италия & 86{,}19 & $+16{,}5$ & 22{,}20 & $-4{,}5$ & $+2{,}08$ \\
8 & Саудовская Аравия & 74{,}35 & $+5{,}6$ & 21{,}18 & $+9{,}9$ & $-8{,}24$ \\
9 & Нидерланды & 67{,}60 & $+20{,}6$ & 19{,}82 & $+17{,}6$ & $-0{,}61$ \\
10 & Бельгия & 65{,}92 & $+27{,}0$ & 23{,}83 & $+7{,}0$ & $+7{,}39$ \\
\dots & & & & & & \\
24 & Респ.\ Корея* (!) & 19{,}44 & $+45{,}6$ & 2{,}23 & $+51{,}2$ & $+15{,}85$ \\
\bottomrule
\end{tabular}

\noindent Источник: GTAIC (2026), табл.~51--52 (стр.~64--99), табл.~45 (стр.~55--59).
Периоды: LTM по странам (стр.~12); доля и CAGR -- за полный 2025 г. «(!)» -- официальное
предупреждение отчёта о вероятной недостоверности тоннажа. {\footnotesize У Бельгии CAGR
в табл.~45 отчёта равен $+7{,}39\,\%$, в сводном резюме -- $+10{,}56\,\%$; расхождение
разделов отчёта не устранено.}
\end{table}

\textbf{США -- треть мирового рынка и главный драйвер роста.} LTM-импорт 736{,}33 млн долл.
($+4{,}1\,\%$) и 264{,}11 тыс.~т ($+1{,}3\,\%$): 29{,}83\,\% стоимости 2025 года и около 33\,\%
тоннажа, CAGR за пять лет $+9{,}36\,\%$ -- единственный устойчивый рост среди большой пятёрки.
Структура поставок сдвигается: Аргентина нарастила долю до 25{,}4\,\% (годом ранее 20{,}0\,\%),
Индия сократила поставки на 27{,}4\,\% (доля 19{,}0\,\% против 27{,}2\,\%), Бразилия -- на
15{,}0\,\%; Новая Зеландия выросла на 46{,}3\,\% до 12{,}7\,\%. Отчёт связывает перелом
с антидемпинговыми пошлинами США с июня 2022 г.\ на сырой мёд из Аргентины, Бразилии, Индии
и Китая (стр.~213); мы цитируем это как контекст отчёта, а не как проверенную причину.

\textbf{Германия} растёт в деньгах ($+10{,}3\,\%$ LTM) при стагнации тоннажа ($-0{,}7\,\%$):
рост стоимости обеспечен ценами; крупнейший поставщик -- Аргентина (18{,}9\,\%), заметна
Украина (13{,}1\,\%, $+41{,}9\,\%$).

\textbf{Япония -- рисковый рынок по классификации отчёта.} Это единственный крупный рынок
с падением по обоим измерителям: LTM 142{,}32 млн долл. ($-3{,}4\,\%$) и 48{,}32 тыс.~т
($-5{,}7\,\%$), пятилетний CAGR $-2{,}84\,\%$, самый низкий прогнозный рост среди 40 стран.
Отчёт объясняет это ослаблением иены со 140{,}5 до 151{,}5 иен за доллар. Поставки
концентрированы: Китай -- 42{,}5\,\% рынка ($-6{,}3\,\%$), Новая Зеландия -- 18{,}5\,\%,
Канада -- 10{,}9\,\% (стр.~195). Для нашей команды это согласуется с наблюдением ДЗ-5:
импорт Японии в долларах сокращался и в 2021--2024 гг.

\textbf{Республика Корея* -- быстрый, но статистически ненадёжный рынок.} LTM 19{,}44 млн
долл. ($+45{,}6\,\%$) при тоннаже всего 2{,}23 тыс.~т: средняя цена-прокси 8\,704 долл./т --
примерно втрое выше средней по рынку (2\,983 долл./т). Отчёт сам помечает корейский тоннаж
флагом достоверности (аномалия CIF) и строит импорт на зеркальной реконструкции с непокрытыми
5{,}28\,\% поставок. Поставщики: Новая Зеландия -- 10{,}96 млн долл., США -- 4{,}45 млн
(стр.~204). Рост в долларах поэтому подтверждается, а «взрывной» рост в тоннах -- нет. Это
тот случай, когда вывод приходится ограничивать измерителем.

\section*{4. Поставщики: цена и политика определяют состав участников}

Топ поставщиков LTM приведён в табл.~4. Рынок высоко концентрирован в физическом измерении:
тройка «Китай -- Аргентина -- Индия» даёт 47{,}0\,\% тоннажа, но лишь около 29\,\% стоимости --
из-за ценового разрыва от 1\,438 долл./т у Китая до 20\,810 долл./т у Новой Зеландии.

\begin{table}[h]
\centering
\small
\caption{Ведущие поставщики мёда в страны топ-40, LTM}
\begin{tabular}{lrrr}
\toprule
Поставщик & млн долл. & Доля, \% & Цена, долл./т \\
\midrule
Аргентина & 293{,}76 & 12{,}02 & 2\,578 \\
Новая Зеландия & 288{,}30 & 11{,}79 & 20\,810 \\
Китай & 231{,}05 & 9{,}45 & 1\,438 \\
Индия & 184{,}61 & 7{,}55 & 1\,783 \\
Украина & 143{,}85 & 5{,}88 & 3\,088 \\
Германия & 126{,}14 & 5{,}16 & 4\,185 \\
\bottomrule
\end{tabular}

\noindent Источник: GTAIC (2026), табл.~54 (стр.~110--124), табл.~59 (стр.~125--140),
табл.~64 (стр.~141). Доли -- от суммарного LTM-импорта 40 стран (около 2{,}44 млрд долл.).
\end{table}

Крупнейшие сдвиги LTM: Аргентина добавила 81{,}06 млн долл.\ и вышла на первое место среди
поставщиков; Индия потеряла 42{,}95 млн долл., в том числе 52{,}7 млн на рынке США и 5{,}03 млн
в ОАЭ* (частично компенсировано ростом на других рынках); Украина развернулась с американского рынка ($-8{,}57$ тыс.~т) на
европейские -- приросты во Францию ($+12{,}1$ млн), Германию ($+10{,}2$ млн) и Польшу
($+6{,}22$ млн). Новая Зеландия занимает лишь 1{,}72\,\% тоннажа, но 11{,}79\,\% стоимости --
всё из-за мануки. В США новозеландский мёд стоит 18{,}7 тыс.\ долл./т при средней цене рынка
2{,}79 тыс., то есть в 6{,}7 раза дороже среднего; во Франции разрыв двенадцатикратный
(45{,}8 тыс.\ против 3{,}85 тыс.\ долл./т). Региональный сдвиг -- иберийский всплеск.
Испания нарастила тоннаж на $47{,}6\,\%$ (крупнейший абсолютный прирост, $+34{,}51$ млн
долл.), Португалия -- на $92{,}9\,\%$. Политические меры ломают отдельные рынки: Словакия потеряла 41{,}9\,\%
тоннажа после эмбарго на украинский мёд с января 2024 г.\ и «экстренного тормоза» ЕС --
Регламента (ЕС) 2024/2166 (стр.~209). По зеркальным данным GACC, поставки России в Китай
сократились на 54{,}5\,\%, до 0{,}79 млн долл.\ LTM; наша панель BACI даёт 2{,}0 млн долл.\
за 2023 г. Расхождение источника, периода и зеркальной методики требует отдельной проверки
и на выбор целевых рынков не влияет.

\section*{5. Сопоставление с расчётами команды}

Наши расчёты выполнены в долларах по OEC/CEPII BACI (HS6 040900, 2018--2024)~\cite{oecbaci};
отчёт добавляет физические объёмы и цены. Сопоставление по трём целевым рынкам -- в табл.~5.

\begin{table}[h]
\centering
\small
\caption{Целевые рынки команды: команда (BACI, 2024) против отчёта (LTM)}
\begin{tabular}{lrllp{4.6cm}}
\toprule
Рынок & BACI 2024, млн долл. & Отчёт, млн долл. & Отчёт, тыс.~т & Цена и примечание \\
\midrule
США & 621{,}2 & 736{,}33 & 264{,}11 & 2\,790 долл./т; согласуется \\
Япония & 132{,}2 & 142{,}32 & 48{,}32 & 2\,950 долл./т; согласуется \\
Южная Корея* & 16{,}5 & 19{,}44 & 2{,}23 (!) & 8\,704 долл./т; тоннаж под вопросом \\
\bottomrule
\end{tabular}

\noindent Источники: расчёты команды по OEC/CEPII BACI (ДЗ-5, дата обращения 19.10.2026
и 04.11.2026); GTAIC (2026), табл.~51--53. Периоды не совпадают (календарный 2024 г.\ против
LTM-окон стр.~12), поэтому сопоставимы порядки величин и направления, а не точные уровни.
\end{table}

Качественно данные согласуются. США -- крупнейший импортёр в обоих источниках; наша оценка
NEX $-0{,}937$ и отчётное доминирование импорта сходятся. Япония в обоих источниках
сокращается: $-21{,}6\,\%$ стоимости за 2021--2024 гг.\ у нас против $-3{,}4\,\%$ LTM и CAGR
$-2{,}84\,\%$ в отчёте. Корея в обоих растёт ($+9{,}4\,\%$ у нас против $+45{,}6\,\%$ LTM --
разница частично объясима несинхронными окнами).

Расхождение измерителей проиллюстрируем одним примером: по стоимости Япония третья
(142{,}32 млн долл.), по тоннажу -- лишь четвёртая, поскольку Великобритания импортирует
62{,}40 тыс.~т по 2\,270 долл./т против японских 2\,950 долл./т. Это прямое подтверждение
нашей оговорки из ДЗ-5: ранжирование рынков по стоимости и по количеству различается, и
долларовая панель BACI, которой мы пользовались, смещает сравнение в сторону дорогих сортов.
Проверка согласованности по Корее: при средней по рынку цене около 3 тыс.\ долл./т импорт
в 19{,}44 млн долл.\ соответствовал бы примерно 6--7 тыс.~т, а не 2{,}23 тыс.~т; это указывает
на недоучёт массы в зеркальных данных, а не на «дорогой» корейский спрос. Такая перекрёстная
проверка -- наш приём, а не вывод отчёта; для окончательного ответа нужен независимый источник
тоннажа.

\section*{6. Критическая оценка отчёта}

\textbf{Сильные стороны.} Во-первых, охват: 96{,}07\,\% мирового импорта в тоннах на стыке трёх
первичных баз (UN Comtrade, Eurostat Comext, GACC) с явным описанием правил стыковки
(стр.~2--9). Во-вторых, культура качества данных: разметка каждого значения по градациям
достоверности веса, официальные предупреждения об аномальных странах и явная формула цен-прокси
- редкость для рыночных отчётов. В-третьих, двустороннее представление (доллары и тонны)
с ценами и единообразные профили всех 40 стран (стр.~177--216), позволяющие проверять любые
сводные цифры по первичным таблицам. В-четвёртых, честные самоограничения: матрица
арбитражных возможностей прямо сопровождается признанием, что она не учитывает пошлины,
расстояния и доступность рынков (стр.~174--175).

\textbf{Ограничения.} (1)~Несинхронные LTM-окна с разбросом до пяти месяцев делают «агрегат
LTM» и межстрановые сравнения формально некорректными. (2)~«Прогнозируемый рост» -- линейная
экстраполяция 24-месячного тренда, а не прогноз: экстремальные значения (Литва $+150\,\%$,
Марокко $+96{,}6\,\%$) -- артефакты малой базы. (3)~Рейтинг привлекательности -- среднее семи
рангов с равными весами без нормировки на размер рынка: первое место Франции (8{,}43 балла)
при 4{,}8\,\% рынка и «худшие» позиции Хорватии, Литвы и Кувейта при положительном росте
показывают, что методика даёт инверсии. (4)~Арбитражная матрица игнорирует пошлины и режимы
доступа -- а именно они, как показывает случай Индии и США, определяют состав участников.
(5)~Обнаруживаются мелкие внутренние рассинхронизации разделов (Бельгия: CAGR $+7{,}39\,\%$
в табл.~45 против $+10{,}56\,\%$ в резюме; Австрия: 25{,}39 против 25{,}61 млн долл.;
Ирландия: 22{,}34 против 22{,}48 млн долл.). (6)~Тоннаж частично модельный: у пяти стран --
зеркальная реконструкция, у Китая с 2025 г.\ -- расчётный вес, у Израиля заявленный вес
отсутствует полностью. (7)~Цены-прокси отражают сортовую структуру поставок (манука), поэтому
«самые дорогие рынки» -- не синоним премиального спроса.

Общая оценка: отчёт корректно описывает \emph{потоки} и аккуратно помечает \emph{качество
данных}, но причины -- цены, курсы, пошлины, урожаи -- упоминает лишь контекстно, без оценки
их влияния; его цифры следует цитировать с заявленной точностью (2--3 значащих цифры) и
предпочтительно по календарным годам, а не по LTM.

\section*{7. Значение для исследования команды}

Отчёт восполняет наш главный информационный пробел: до сих пор все расчёты команды были
в долларах, и оговорка «рост стоимости не доказывает рост количества» оставалась
непроверяемой. Теперь по каждому целевому рынку доступны тонны и цены-прокси, что позволяет
(а) перепроверить выводы ДЗ-5 в физическом измерении, (б) оценить ценовой сегмент, в который
мы потенциально поставляем, и (в) учесть конкурентный фон: вытеснение Индии с рынка США
на фоне антидемпинговых пошлин, доминирование Китая (42{,}5\,\%) и валютный риск в Японии,
быстрый, но статистически ненадёжный рост в Корее. Крупнейший в отчёте «разрыв спроса --
предложения» приходится на США (14{,}96 млн долл./год), что согласуется с нашим отбором США
как базового рынка в ДЗ-8--ДЗ-9. Важный отрицательный результат: рейтинг привлекательности
GTAIC и его разрывы не учитывают тарифы, ветеринарные требования и логистику -- то, что
составляет предмет нашего ДЗ-12, поэтому рейтинг используем как гипотезу, а не как основание
выбора. Следующий шаг -- сверить корейский тоннаж по независимому источнику (национальная
статистика Кореи или зеркало ЕС) и лишь затем фиксировать корейский спрос в тоннах; параллельно
- сопоставить ценовые градации отчёта с нашим продуктовым позиционированием.

\section*{Источники}
\begin{thebibliography}{9}
\bibitem{bradbear} Bradbear, N. (2004). \textit{Beekeeping and Sustainable Livelihoods}.
FAO Diversification Booklet No.~1. Rome: FAO.
\url{https://www.fao.org/4/y5110e/y5110e02.htm} (дата обращения: 19.10.2026).
Роль: товарная характеристика мёда и медоносных ресурсов для факторной гипотезы ДЗ-5.
\bibitem{gtaic2026} EXPORT HUNTER, UAB (2026). \textit{Top-40 Importing Countries World
040900. Volume in Tons} [Отчёт платформы GTAIC]. Вильнюс, 217~с. Сгенерирован 24.11.2026.
\url{https://www.gtaic.ai} (дата обращения: 25.11.2026).
Роль: предмет обзора; масштаб и динамика спроса, физические объёмы, цены-прокси, поставщики
и торговые меры.
\bibitem{oecbaci} Observatory of Economic Complexity / CEPII (2026). \textit{BACI International
Trade Database}, куб \textit{trade\_i\_baci\_a\_17}, HS6 REV.~2017, код 040900, 2018--2024,
долл. США. \url{https://api-v2.oec.world/tesseract} (даты обращения: 19.10.2026, 04.11.2026).
Роль: собственные расчёты команды (ДЗ-5, ДЗ-8, ДЗ-9) -- долларовая панель импорта и экспорта.
\bibitem{wdi} World Bank (2026). \textit{World Development Indicators}: AG.LND.ARBL.HA.PC
(пашня, га на человека, 2024), NY.GDP.MKTP.CD, SP.POP.TOTL, PA.NUS.FCRF.
\url{https://api.worldbank.org/v2/} (обновление 08.10.2026; дата обращения: 19.10.2026).
Роль: факторная обеспеченность и переменные панельных моделей команды.
\end{thebibliography}

\end{document}
